\documentclass[UTF8,11pt]{ctexart}
\usepackage[a4paper,margin=22mm]{geometry}
\usepackage{graphicx,enumitem,amssymb,tabularx}
\usepackage[normalem]{ulem}
\usepackage{xeCJKfntef}
\setlength{\parindent}{0pt}
\setlength{\parskip}{.6em}
\sloppy
\begin{document}
\section*{2015年6月英语六级真题(第3套)}
\subsection*{2015 年6 月大学英语六级考试真题(三)}


\subsection*{\textbf{Part I} \textbf{Writing} \textbf{(30 minutes)}}


Directions：\textit{For this part},\textit{ you are allowed 30 minutes to write an essay commenting on the saying ulf you cannot do great}


\subsection*{\textit{things},\textit{ do small things in a great way.99 You can cite examples to illustrate your point of view. You should write at least} \textbf{\textit{150}}}


\subsection*{\textit{words but no more than}\textbf{\textit{ 200}}\textit{ words.}}


注意：此部分试题请在答题卡1 上作答。


\textbf{Part}II \textbf{Listening Comprehension} \textbf{(30 minutes)}


\subsection*{\textbf{Section A}}


Directions：\textit{In this section},\textit{ you will hear 8 short conversations and 2 long conversations. At the end of each conversation,} \textit{one or more questions will be asked about what was said. Both the conversation and the questions will be spoken only once.} \textit{After each question there will be a pause}.\textit{ During the pause},\textit{ you must read the four choices marked A) 9 B),} C)\textit{ and D), and}


\subsection*{\textit{decide which is the best answer. Then mark the corresponding letter on}\textbf{\textit{ Answer Sheet 1}}\textit{ with a single line through the}}


centre. 注意：此部分试题请在答题卡1 上作答。


1.


\begin{itemize}[leftmargin=2em,itemsep=.2em,topsep=.4em]\item[A.] The man might be able to play in the World Cup.
\end{itemize}

\begin{itemize}[leftmargin=2em,itemsep=.2em,topsep=.4em]\item[B.] The man’s football career seems to be at an end.
\item[C.] The man was operated on a few weeks ago.
\item[D.] The man is a fan of world-famous football players.
\end{itemize}

2.


\begin{itemize}[leftmargin=2em,itemsep=.2em,topsep=.4em]\item[A.] Work out a plan to tighten his budget.
\item[C.] Apply for a senior position in the restaurant.
\end{itemize}

\begin{itemize}[leftmargin=2em,itemsep=.2em,topsep=.4em]\item[B.] Find out the opening hours of the cafeteria.
\item[D.] Solve his problem by doing a part-time job.
\end{itemize}

3.


\begin{itemize}[leftmargin=2em,itemsep=.2em,topsep=.4em]\item[A.] A financial burden.
\item[C.] A real nuisance.
\end{itemize}

\begin{itemize}[leftmargin=2em,itemsep=.2em,topsep=.4em]\item[B.] A good companion.
\item[D.] A well-trained pet.
\end{itemize}

4.


\begin{itemize}[leftmargin=2em,itemsep=.2em,topsep=.4em]\item[A.] The errors will be corrected soon.
\item[C.] The computing system is too complex.
\end{itemize}

\begin{itemize}[leftmargin=2em,itemsep=.2em,topsep=.4em]\item[B.] The woman was mistaken herself.
\item[D.] He has called the woman several times.
\end{itemize}

5.


\begin{itemize}[leftmargin=2em,itemsep=.2em,topsep=.4em]\item[A.] He needs help to retrieve his files.
\item[C.] He needs some time to polish his paper.
\end{itemize}

\begin{itemize}[leftmargin=2em,itemsep=.2em,topsep=.4em]\item[B.] He has to type his paper once more.
\item[D.] He will be away for a two-week conference.
\end{itemize}

6.


\begin{itemize}[leftmargin=2em,itemsep=.2em,topsep=.4em]\item[A.] They might have to change their plan.
\item[C.] He has a heavier workload than the woman.
\end{itemize}

\begin{itemize}[leftmargin=2em,itemsep=.2em,topsep=.4em]\item[B.] He has got everything set for their trip.
\item[D.] They could stay in the mountains until June 8.
\end{itemize}

7.


\begin{itemize}[leftmargin=2em,itemsep=.2em,topsep=.4em]\item[A.] They have to wait a month to apply for a student loan.
\end{itemize}

\begin{itemize}[leftmargin=2em,itemsep=.2em,topsep=.4em]\item[B.] They can find the application forms in the brochure.
\item[C.] They are not eligible for a student loan.
\item[D.] They are not late for a loan application.
\end{itemize}

8.


\begin{itemize}[leftmargin=2em,itemsep=.2em,topsep=.4em]\item[A.] New laws are yet to be made to reduce pollutant release.
\end{itemize}

\begin{itemize}[leftmargin=2em,itemsep=.2em,topsep=.4em]\item[B.] Pollution has attracted little attention from the public.
\item[C.] The quality of air will surely change for the better.
\item[D.] It’ll take years to bring air pollution under control.
\end{itemize}

Questions 9 to 12 are based on the conversation you have just heard


9.


\begin{itemize}[leftmargin=2em,itemsep=.2em,topsep=.4em]\item[A.] Enormous size of its stores.
\item[C.] Its appealing surroundings.
\end{itemize}

\begin{itemize}[leftmargin=2em,itemsep=.2em,topsep=.4em]\item[B.] Numerous varieties of food.
\item[D.] Its rich and colorful history.
\end{itemize}

10.


\begin{itemize}[leftmargin=2em,itemsep=.2em,topsep=.4em]\item[A.] An ancient building.
\item[C.] An Egyptian museum.
\end{itemize}

\begin{itemize}[leftmargin=2em,itemsep=.2em,topsep=.4em]\item[B.] A world of antiques.
\item[D.] An Egyptian memorial.
\end{itemize}

11.


\begin{itemize}[leftmargin=2em,itemsep=.2em,topsep=.4em]\item[A.] Its power bill reaches £ 9 million a year.
\item[C.] It supplies power to a nearby town.
\end{itemize}

\begin{itemize}[leftmargin=2em,itemsep=.2em,topsep=.4em]\item[B.] It sells thousands of light bulbs a day.
\item[D.] It generates 70\% of the electricity it uses.
\end{itemize}

12.


\begin{itemize}[leftmargin=2em,itemsep=.2em,topsep=.4em]\item[A.] 11 500.
\item[B.] 30 000.
\item[C.] 250 000.
\item[D.] 300 000.
\end{itemize}

Questions 13 to 15 are based on the conversation you have just heard.


13.


\begin{itemize}[leftmargin=2em,itemsep=.2em,topsep=.4em]\item[A.] Transferring to another department.
\item[C.] Thinking about doing a different job.
\end{itemize}

\begin{itemize}[leftmargin=2em,itemsep=.2em,topsep=.4em]\item[B.] Studying accounting at a university.
\item[D.] Making preparations for her wedding.
\end{itemize}

14.


\begin{itemize}[leftmargin=2em,itemsep=.2em,topsep=.4em]\item[A.] She has finally got a promotion and a pay raise.
\end{itemize}

\begin{itemize}[leftmargin=2em,itemsep=.2em,topsep=.4em]\item[B.] She has got a satisfactory job in another company.
\item[C.] She could at last leave the accounting department.
\item[D.] She managed to keep her position in the company.
\end{itemize}

15.


\begin{itemize}[leftmargin=2em,itemsep=.2em,topsep=.4em]\item[A.] He and Andrea have proved to be a perfect match.
\end{itemize}

\begin{itemize}[leftmargin=2em,itemsep=.2em,topsep=.4em]\item[B.] He changed his mind about marriage unexpectedly.
\item[C.] He declared that he would remain single all his life.
\item[D.] He would marry Andrea even without meeting her.
\end{itemize}

\subsection*{\textbf{Section B}}


Directions：\textit{In this section you will hear 3 short passages. At the end of each passage, you will hear some questions.Both the}


\subsection*{\textit{passage and the questions will be spoken only once. After you hear a question}，\textit{you must choose the best answer from the}}


\subsection*{\textit{four choices marked A} ),\textit{ B}), C)\textit{ and D). Then mark the corresponding letter on} \textbf{\textit{AnswerSheet1}} \textit{with a single line through the}}


\textit{centre.} 注意：此部分试题请在答题卡1 上作答。


\subsection*{Passage One}


Questions 16 to 19 are based on the passage you have just heard.


16.


\begin{itemize}[leftmargin=2em,itemsep=.2em,topsep=.4em]\item[A.] They are motorcycles designated for water sports.
\end{itemize}

\begin{itemize}[leftmargin=2em,itemsep=.2em,topsep=.4em]\item[B.] They are speedy boats restricted in narrow waterways.
\item[C.] They are becoming an efficient form of water transportation.
\item[D.] They are getting more popular as a means of water recreation.
\end{itemize}

17.


\begin{itemize}[leftmargin=2em,itemsep=.2em,topsep=.4em]\item[A.] Water scooter operators’ lack of experience.
\item[C.] Overloading of small boats and other craft.
\end{itemize}

\begin{itemize}[leftmargin=2em,itemsep=.2em,topsep=.4em]\item[B.] Vacationers’ disregard of water safety rules.
\item[D.] Carelessness of people boating along the shore.
\end{itemize}

18.


\begin{itemize}[leftmargin=2em,itemsep=.2em,topsep=.4em]\item[A.] They scare whales to death.
\item[C.] They discharge toxic emissions.
\end{itemize}

\begin{itemize}[leftmargin=2em,itemsep=.2em,topsep=.4em]\item[B.] They produce too much noise.
\item[D.] They endanger lots of water life.
\end{itemize}

19.


\begin{itemize}[leftmargin=2em,itemsep=.2em,topsep=.4em]\item[A.] Expand operating areas.
\item[B.] Limit the use of water scooters.
\end{itemize}

\begin{itemize}[leftmargin=2em,itemsep=.2em,topsep=.4em]\item[C.] Restrict operating hours.
\item[D.] Enforce necessary regulations.
\end{itemize}

\subsection*{Passage Two}


Questions 20 to 22 are based on the passage you have just heard.


20.


\begin{itemize}[leftmargin=2em,itemsep=.2em,topsep=.4em]\item[A.] They are stable.
\item[B.] They are close.
\item[C.] They are strained.
\item[D.] They are changing.
\end{itemize}

21.


\begin{itemize}[leftmargin=2em,itemsep=.2em,topsep=.4em]\item[A.] They are fully occupied with their own business.
\end{itemize}

\begin{itemize}[leftmargin=2em,itemsep=.2em,topsep=.4em]\item[B.] Not many of them stay in the same place for long.
\item[C.] Not many of them can win trust from their neighbors.
\item[D.] They attach less importance to interpersonal relations.
\end{itemize}

22.


\begin{itemize}[leftmargin=2em,itemsep=.2em,topsep=.4em]\item[A.] Count on each other for help.
\item[C.] Keep a friendly distance.
\end{itemize}

\begin{itemize}[leftmargin=2em,itemsep=.2em,topsep=.4em]\item[B.] Give each other a cold shoulder.
\item[D.] Build a fence between them.
\end{itemize}

\subsection*{Passage Three}


Questions 23 to 25 are based on the passage you have just heard.


23.


\begin{itemize}[leftmargin=2em,itemsep=.2em,topsep=.4em]\item[A.] It may produce an increasing number of idle youngsters.
\end{itemize}

\begin{itemize}[leftmargin=2em,itemsep=.2em,topsep=.4em]\item[B.] It may affect the quality of higher education in America.
\item[C.] It may cause many schools to go out of operation.
\item[D.] It may lead to a lack of properly educated workers.
\end{itemize}

24.


\begin{itemize}[leftmargin=2em,itemsep=.2em,topsep=.4em]\item[A.] It is less serious in cities than in rural areas.
\end{itemize}

\begin{itemize}[leftmargin=2em,itemsep=.2em,topsep=.4em]\item[B.] It affects both junior and senior high schools.
\item[C.] It results from a worsening economic climate.
\item[D.] It is a new challenge facing American educators.
\end{itemize}

25.


\begin{itemize}[leftmargin=2em,itemsep=.2em,topsep=.4em]\item[A.] Allowing them to choose their favorite teachers.
\end{itemize}

\begin{itemize}[leftmargin=2em,itemsep=.2em,topsep=.4em]\item[B.] Creating a more relaxed learning environment.
\item[C.] Rewarding excellent academic performance.
\item[D.] Helping them to develop better study habits.
\end{itemize}

\subsection*{\textbf{Section C}}


Directions：\textit{In this section}，\textit{you will hear a passage three times}．\textit{When the passage is read for the first time}，\textit{you should} \textit{listen carefully for its general idea}．\textit{When the passage is read for the second time}，\textit{you are required to fill in the blanks with} \textit{the exact words you have just heard}．\textit{Finally}，\textit{when the passage is read for the third time}，\textit{you should check what you have} \textit{written}． 注意：此部分试题请在答题卡1 上作答。


I'm interested in the criminal justice system of our country．It seems to me that something has to be done，if we’re to\uline{ 26} as a country．I certainly don’t know what the answers to our problems are．Things certainly get\uline{ 27 }in a hurry when you get into them，but l wonder if something couldn’t be done to deal with some of these problems．One thing I'm concerned about is our practice of putting\uline{ 28} in jail who haven’t harmed anyone．Why not work out some system whereby they can pay back the debts they owe society instead of \uline{29} another debt by going to prison and，of course，coming‘\uline{ 30 }hardened criminals．I'm also concerned about the short prison sentences people are\uline{ 31 }serious crimes．Of course one alternative to this is to\uline{ 32 }capital punishment，but I'm not sure l would be for that．I'm not sure it’s right to take an eye for an eye．The alternative to capital punishment is longer sentences。but they would certainly cost the tax payers much money．I also think we must do something about the insanity\uline{ 33 }．In my opinion，anyone who takes another person’s life \uline{34 }is insane，however，that does not mean that the person isn’t guilty of the crime，or that he shouldn’t pay society the debt he owes．It’s sad，of course，that a person may have to spend the rest of his life，or a large part of it in prison for acts that he\uline{ 35 }while not in full control of his mind．


\subsection*{\textbf{Part} III \textbf{Reading Comprehension} \textbf{(40 minutes)}}


\subsection*{\textbf{Section A}}


Directions ：\textit{In this section},\textit{ there is a passage with ten blanks. You are required to select one word for each blank from a list} \textit{of choices given in a word bank following the passage. Read the passage through carefully before making your choices.}


\subsection*{\textit{Each choice in the bank is identified by a letter. Please mark the corresponding letter for each item on}\textbf{\textit{ Answer}} \textbf{\textit{Sheet}} \textbf{\textit{2}}\textit{ with}}


\textit{a single line through the centre. You may not use any of the words in the bank more than once}.


Questions 36 to 45 are based on the following passage.


Travel websites have been around since the 1990s, when Expedia, Travelocity, and other holiday booking sites were launched, allowing travelers to compare flight and hotel prices with the click of a mouse. With information no longer \uline{36} by travel agents or hidden in business networks, the travel industry was revolutionized, as greater transparency helped \uline{37} prices.


Today, the industry is going through a new revolution—this time transforming service quality. Online rating platforms —\uline{ 38} in hotels, restaurants, apartments, and taxis —allow travelers to exchange reviews and experiences for all to see.


Hospitality businesses are now ranked, analyzed, and compared not by industry \uline{39 }，but by the very people for whom the service is intended—the customer. This has \uline{40} a new relationship between buyer and seller. Customers have always voted with their feet；they can now explain their decision to anyone who is interested. As a result, businesses are much more \uline{41 }，often in very specific ways, which creates powerful \uline{42} to improve service.


Although some readers might not care for gossipy reports of unfriendly\textit{ bellboys}(行李员）in Berlin or malfunctioning hotel hairdryers in Houston, the true power of online reviews lies not just in the individual stories, but in the websites’ \uline{43} to aggregate a large volume of ratings.


The impact cannot be \uline{44} Businesses that attract top ratings can enjoy rapid growth, as new customers are attracted by good reviews and \uline{45} provide yet more positive feedback. So great is the influence of online ratings that many companies now hire digital reputation managers to ensure a favorable online identity.


注意：此部分试题请在答题卡2 上作答。


\begin{itemize}[leftmargin=2em,itemsep=.2em,topsep=.4em]\item[A.] accountable
\item[B.] capacity
\item[C.] controlled
\item[D.] entail
\item[E.] forged
\item[F.] incentives
\item[G.] occasionally
\item[H.] overstated
\item[I.] persisting
\item[J.] pessimistic
\item[K.] professionals
\item[L.] slash
\item[M.] specializing
\item[N.] spectators
\item[O.] subsequently
\end{itemize}

\subsection*{\textbf{Section B}}


Directions：\textit{In this section, you are going to read a passage with ten statements attached to it. Each statement contains} \textit{information given in one of the paragraphs. Identify the paragraph from which the information is derived. You may choose} \textit{a paragraph more than once. Each paragraph is marked with a letter . Answer the questions by marking the corresponding}


\subsection*{\textit{letter on}\textbf{\textit{ Answer}} \textbf{\textit{Sheet}} \textbf{\textit{2}}.}


\textbf{Plastic Surgery} A better credit card is the solution to ever larger hack attacks .


A) A thin magnetic strip (magstripe) is all that stands between your credit-card information and the bad guys. And they’e been working hard to break in. That’s why 2014 is shaping up as a major showdown：banks, law enforcement and technology companies are all trying to stop a network of hackers who are succeeding in stealing account numbers, names, email addresses and other crucial data used in identity theft. More than 100 million accounts at Target, Neiman Marcus and Michaels stores were affected in some way during the most recent attacks, starting last November.


B) \textit{Swipe} (刷卡) is the operative word：cards are increasingly vulnerable to attacks when you make purchases in a store. In several recent incidents, hackers have been able to obtain massive information of credit-\textit{debit-}(借记）or prepaid-card numbers using malware, i. e. malicious software, inserted secretly into the retailers’ point-of-sale system—the checkout registers. Hackers then sold the data to a second group of criminals operating in shadowy corners of the web. Not long after, the stolen data was showing up on fake cards and being used for online purchases.


C) The solution could cost as little as \$2 extra for every piece of plastic issued. The fix is a security technology used heavily outside the US. While American credit cards use the 40-year-old magstripe technology to process transactions, much of the rest of the world uses smarter cards with a technology called EMV (short for Europay, MasterCard, Visa) that employs a chip embedded in the card plus a customer PIN (personal identification number) to\textit{ authenticate}(验证)every transaction on the spot. If a purchaser fails to punch in the correct PIN at the checkout, the transaction gets rejected. (Online purchases can be made by setting up a separate transaction code.)


D) Why haven’t big banks adopted the more secure technology? When it comes to mailing out new credit cards, it’s all about relative costs, says David Robertson, who runs the\textit{ Nilson Report,} an industry newsletter. “The cost of the card, putting the sticker on it, coding the account number and expiration date,\textit{ embossing}(凸印）it, the small envelope—all put together, you’re in the dollar range. ” A chip- and-PIN card currently costs closer to \$3, says Robertson, because of the price of chips. (Once large issuers convert together, the chip costs should drop.)


E) Multiply \$3 by the more than 5 billion magstripe credit and prepaid cards in circulation in the US. Then consider that there’s an estimated \$12.4 billion in card fraud on a global basis, says Robertson. With 44\% of that in the US, American credit-card fraud amounts to about \$5.5 billion annually. Card issuers have so far calculated that absorbing the liability for even big hacks like the Target one is still cheaper than replacing all that plastic.


F) That leaves American retailers pretty much alone the world over in relying on magstripe technology to charge purchases —and leaves consumers vulnerable. Each magstripe has three tracks of information, explains payments security expert Jeremy Gumbley, the chief technology officer of Credit Call, an electronic-payments company. The first and third are used by the bank or card issuer. Your vital account information lives on the second track, which hackers try to capture. “Malware is scanning through the memory in real time and looking for data，” he says. “It creates a text file that gets stolen. ”


G) Chip-and-PIN cards, by contrast, make fake cards or skimming impossible because the information that gets scanned is \textit{encrypted} (加密). The historical reason the US has stuck with magstripe, ironically enough, is once superior technology. Our cheap, ultra-reliable wired networks made credit- card authentication over the phone frictionless. In France, card companies created EMV in part because the telephone monopoly was so maddeningly inefficient and expensive. The EMV solution allowed transactions to be verified locally and securely.


H) Some big banks, like Wells Fargo, are now offering to convert your magstripe card to a chip-and-PIN model. (It’s actually a \textit{hybird} (混合体）that will still have a magstripe, since most US merchants don’t have EMV terminals.) Should you take them up on it? If you travel internationally, the answer is yes.


I) Keep in mind, too, that credit cards typically have better liability protection than debit cards. If someone uses your credit card \textit{fraudulently}（欺诈性的），it’s the issuer or merchant, not you, that takes the hit. Debit cards have different liability limits depending on the bank and the events surrounding any fraud. “If it’s available, the logical thing is to get a chip-and-PIN card from your bank，” says Eric Adamowsky, a co-founder of Credit Cardlnsider. com. “I would use credit cards over debit cards because of liability issues. ’’ Cash still works pretty well too.


J) Retailers and banks stand to benefit from the lower fraud levels of chip-and-PIN cards but have been reluctan for years to invest in the new\textit{ infrastructure} (基袖设施）needed for the technology，especially if consumers don＇t have access to it. It＇ s a chicken-and-egg problem：no one wants to spend the money on upgraded point-of-sale systems that can read the chip cards if shoppers aren’t carrying them—yet there’s little point in consumers’ carrying the fancy plastic if stores aren’t equipped to use them. (An earlier effort by Target to move to chip and PIN never gained progress.) According to Gumbley, there’s a “you-first mentality. The\textit{ logjam} (僵局）has to be broken.”


K) JPMorgan Chase CEO Jamie Dimon recently expressed his willingness to do so, noting that banks and merchants have spent the past decade suing each other over interchange fees —the percentage of the transaction price they keep —rather than deal with the growing hacking problem. Chase offers a chip- enabled card under its own brand and several others for travel-related companies such as British Airways and Ritz-Carlton.


L) The Target and Neiman hacks have also changed the cost calculation：although retailers have been reluctant to spend the \$6.75 billion that Capgemini consultants estimate it will take to convert all their registers to be chip-and-PIN-compatible, the potential liability they now face is dramatically greater. Target has been hit with class actions from hacked consumers. “It’s the ultimate nightmare,” a retail executive from a well-known chain admitted to\textit{ TIME.}


M) The card-payment companies MasterCard and Visa are pushing hard for change. The two firms have warned all parties in the transaction chain 一merchant, network, bank —that if they don’t become EMV-compliant by October 2015, the party that is least compliant will bear the fraud risk.


N) In the meantime, app-equipped smartphones and digital wallets—all of which can use EMV technology—are beginning to make\textit{ inroads}（侵染）on cards and cash. PayPal, for instance, is testing an app that lets you use your mobile phone to pay on the fly at local merchants—without surrendering any card information to them. And further down the road is biometric authentication, which could be encrypted with, say, a fingerprint.


O) Credit and debit cards, though, are going to be with us for the foreseeable future, and so are hackers, if we stick with magstripe technology. “It seems crazy to me,” says Gumbley, who is English, “that a cutting-edge-technology country is depending on a 40-year-old technology.That’s why it may be up to consumers to move the needle on chip and PIN. Says Robertson:” When you get the consumer into a position of worry and inconvenience, that’s where the rubber hits the road. ’’


注意：此部分试题请在答题卡2 上作答。


46. It is best to use an EMV card for international travel.


47. Personal information on credit and debit cards is increasingly vulnerable to hacking.


48. The French card companies adopted EMV technology partly because of inefficient telephone service.


49. While many countries use the smarter EMV cards, the US still clings to its old magstripe technology.


50. Attempts are being made to prevent hackers from carrying out identity theft.


51. Credit cards are much safer to use than debit cards.


52. Big banks have been reluctant to switch to more secure technology because of the higher costs involved.


53. The potential liability for retailers using magstripe is far more costly than upgrading their registers.


54. The use of magstripe cards by American retailers leaves consumers exposed to the risks of losing account information.


55. Consumers will be a driving force behind the conversion from magstripe to EMV technology.


\subsection*{\textbf{Section C}}


Directions：\textit{There are 2 passages in this section. Each passage is followed by some questions or unfinished statements • For} \textit{each of them there are four choices marked A) , B),} C)\textit{ and D). You should decide on the best choice and mark the}


\subsection*{\textit{corresponding letter on}\textbf{\textit{ Answer}} \textbf{\textit{Sheet}} \textbf{\textit{2}}\textit{ with a single line through the centre.}}


\subsection*{\textbf{Passage One}}


\textbf{Questions 56 to 60 are based on the following passage.}


The report from the Bureau of Labor Statistics was just as gloomy as anticipated. Unemployment in January jumped to a 16-year high of 7. 6 percent, as 598 000 jobs were slashed from US payrolls in the worst single-month decline since December, 1974. With 1. 8 million jobs lost in the last three months, there is urgent desire to boost the economy as quickly as possible. But Washington would do well to take a deep breath before reacting to the grim numbers.


Collectively, we rely on the unemployment figures and other statistics to frame our sense of reality. They are a vital part of an array of data that we use to assess if we’re doing well or doing badly, and that in turn shapes government policies and corporate budgets and personal spending decisions. The problem is that the statistics aren’t an objective measure of reality；they are simply a best approximation. Directionally, they capture the trends, but the idea that we know precisely how many are unemployed is a myth. That makes finding a solution all the more difficult.


First, there is the way the data is assembled. The official unemployment rate is the product of a telephone survey of about 60 000 homes. There is another survey, sometimes referred to as the “payroll survey” that assesses 400 000 businesses based on their reported payrolls. Both surveys have problems. The payroll survey can easily double-count someone: if you are one person with two jobs, you show up as two workers. The payroll survey also doesn’t capture the number of self-employed, and so says little about how many people are generating an independent income.


The household survey has a larger problem. When asked straightforwardly, people tend to lie or shade the truth when the subject is sex, money or employment. If you get a call and are asked if you’re employed, and you say yes,you’re employed. If you say no, however, it may surprise you to learn that you are only unemployed if you’ve been actively looking for work in the past four weeks；otherwise, you are “marginally attached to the labor force” and not actually unemployed.


The urge to quantify is embedded in our society. But the idea that statisticians can then capture an objective reality isn’t just impossible. It also leads to serious mis judgments. Democrats and Republicans can and will take sides on a number of issues, but a more crucial concern is that both are basing major policy decisions on guesstimates rather than looking at the vast wealth of raw data with a critical eye and an open mind. 注意：此部分试题请在答题卡2 上作答。


56. What do we learn from the first paragraph?


\begin{itemize}[leftmargin=2em,itemsep=.2em,topsep=.4em]\item[A.] The US economic situation is going from bad to worse.
\item[B.] Washington is taking drastic measures to provide more jobs.
\item[C.] The US government is slashing more jobs from its payrolls.
\item[D.] The recent economic crisis has taken the US by surprise.
\end{itemize}

57. What does the author think of the unemployment figures and other statistics?


\begin{itemize}[leftmargin=2em,itemsep=.2em,topsep=.4em]\item[A.] They form a solid basis for policy making.
\item[B.] They represent the current situation.
\item[C.] They signal future economic trends.
\item[D.] They do not fully reflect the reality.
\end{itemize}

58. One problem with the payroll survey is that\par\noindent\rule{.95\linewidth}{.4pt}\par  .


\begin{itemize}[leftmargin=2em,itemsep=.2em,topsep=.4em]\item[A.] it does not include all the businesses
\item[B.] it fails to count in the self-employed
\item[C.] it magnifies the number of the jobless
\item[D.] it does not treat all companies equally
\end{itemize}

59. The household survey can be faulty in that\par\noindent\rule{.95\linewidth}{.4pt}\par  .


\begin{itemize}[leftmargin=2em,itemsep=.2em,topsep=.4em]\item[A.] people tend to lie when talking on the phone
\item[B.] not everybody is willing or ready to respond
\item[C.] some people won’t provide truthful information
\item[D.] the definition of unemployment is too broad
\end{itemize}

60. At the end of the passage,the author suggests that\par\noindent\rule{.95\linewidth}{.4pt}\par  .


\begin{itemize}[leftmargin=2em,itemsep=.2em,topsep=.4em]\item[A.] statisticians improve their data-assembling methods
\end{itemize}

\begin{itemize}[leftmargin=2em,itemsep=.2em,topsep=.4em]\item[B.] decision makers view the statistics with a critical eye
\item[C.] politicians listen more before making policy decisions
\item[D.] Democrats and Republicans cooperate on crucial issues
\end{itemize}

\subsection*{\textbf{Passage Two}}


\textbf{Questions 61 to 65 are based on the following passage.}


At some point in 2008, someone, probably in either Asia or Africa, made the decision to move from the countryside to the city. This nameless person pushed the human race over a historic threshold, for it was in that year that mankind became, for the first time in its history, a predominantly urban species.


It is a trend that shows no sign of slowing.\textit{ Demographers}\textbf{\textit{ (}}人口统计学家）reckon that three-quarters of humanity could be city-dwelling by 2050, with most of the increase coming in the fast-growing towns of Asia and Africa. Migrants to cities are attracted by plentiful jobs, access to hospitals and education, and the ability to escape the boredom of a farmer’s agricultural life. Those factors are more than enough to make up for the\textit{ squalor} (肮脏.), disease and spectacular poverty that those same migrants must often at first endure when they become urban dwellers.


It is the city that inspires the latest book from Peter Smith. His main thesis is that the buzz of urban life, and the opportunities it offers for co-operation and collaboration, is what attracts people to the city, which in turn makes cities into the engines of art, commerce, science and progress. This is hardly revolutionary, but it is presented in a charming format. Mr. Smith has written a breezy guidebook, with a series of short chapters dedicated to specific aspects of urbanity—parks, say, or the various schemes that have been put forward over the years for building the perfect city. The result is a sort of high-quality, unusually rigorous coffee-table book, designed to be dipped into rather than read from beginning to end.


In the chapter on skyscrapers, for example, Mr. Smith touches on construction methods, the revolutionary invention of the automatic lift, the practicalities of living in the sky and the likelihood that, as cities become more crowded, apartment living will become the norm. But there is also time for brief diversions onto bizarre ground, such as a discussion of the skyscraper index (which holds that a boom in skyscraper construction is a foolproof sign of an imminent recession) .


One obvious criticism is that the price of breadth is depth; many of Mr. Smith’s essays raise as many questions as they answer. Although that can indeed be frustrating, this is probably the only way to treat so grand a topic. The city is the building block of civilisation and of almost everything people do；a guidebook to the city is really, therefore, a guidebook to how a large and ever-growing chunk of humanity chooses to live. Mr. Smithes book serves as an excellent introduction to a vast subject, and will suggest plenty of further lines of inquiry. 注意：此部分试题请在答题卡2 上作答。


61. In what way is the year 2008 historic?


\begin{itemize}[leftmargin=2em,itemsep=.2em,topsep=.4em]\item[A.] For the first time in history, urban people outnumbered rural people.
\item[B.] An influential figure decided to move from the countryside to the city.
\item[C.] It is in this year that urbanisation made a start in Asia and Africa.
\item[D.] The population increase in cities reached a new peak in Asia and Africa.
\end{itemize}

62. What does the author say about urbanisation?


\begin{itemize}[leftmargin=2em,itemsep=.2em,topsep=.4em]\item[A.] Its impact is not easy to predict.
\item[B.] Its process will not slow down.
\item[C.] It is a milestone in human progress.
\item[D.] It aggravates the squalor of cities.
\end{itemize}

63. How does the author comment on Peter Smith’s new book?


\begin{itemize}[leftmargin=2em,itemsep=.2em,topsep=.4em]\item[A.] It is but an ordinary coffee-table book.
\item[B.] It is flavoured with humourous stories.
\item[C.] It serves as a guide to arts and commerce.
\item[D.] It is written in a lively and interesting style.
\end{itemize}

64. What does the author say in the chapter on skyscrapers?


\begin{itemize}[leftmargin=2em,itemsep=.2em,topsep=.4em]\item[A.] The automatic lift is indispensable in skyscrapers.
\item[B.] People enjoy living in skyscrapers with a view.
\item[C.] Skyscrapers are a sure sign of a city’s prosperity.
\item[D.] Recession closely follows a skyscraper boom.
\end{itemize}

65. What may be one criticism of Mr. Smith’s book?


\begin{itemize}[leftmargin=2em,itemsep=.2em,topsep=.4em]\item[A.] It does not really touch on anything serious.
\item[B.] It is too long for people to read from cover to cover.
\item[C.] It does not deal with any aspect of city life in depth.
\item[D.] It fails to provide sound advice to city dwellers.
\end{itemize}

\subsection*{\textbf{Part} IV \textbf{Translation} \textbf{(30 minutes)}}


\textbf{Directions}：\textit{For this part, you are allowed 30 minutes to translate a passage from Chinese into English. You should write}


\subsection*{\textit{your answer on}\textbf{\textit{ Answer}} \textbf{\textit{Sheet}} \textbf{\textit{2.}}}


汉朝是中国历史上最重要的朝代之一。汉朝统治期间有很多显著的成就。它最先向其他文化敞开大门，对外贸易兴旺。 汉朝开拓的丝绸之路通向了中西亚乃至罗马。各类艺术一派繁荣，涌现了很多文学、历史、哲学巨著。公元100 年中国第一部字典编撰完成，收入9 000 个字，提供释义并列举不同的写法。其间，科技方面也取得了很大进步，发明了纸张、水钟、日晷（sundials) 以及测量地震的仪器。汉朝历经 400 年，但统治者的腐败最终导致了它的灭亡。 注意:此部分试题请在答题卡2 上作答。

\end{document}